%!TEX program = lualatex
\documentclass{loopmind-doc}

\doclang{es}
\doctype{Resumen de producto}
\doctitle{Makis\newline PRD en una mirada}
\docsubtitle{Qué construye Makis, cómo funciona y qué debe demostrar en el hackathon.}
\docclient{Equipo Maki Acevichado}
\docdate{Septiembre 2026}
\docversion{1.0}

\usetikzlibrary{positioning,arrows.meta}
\setlength{\headheight}{20pt}

\newcommand{\phase}[3]{%
\begin{tcolorbox}[colback=lm-surface,colframe=lm-border,boxrule=.45pt,arc=2mm,
left=4mm,right=4mm,top=3.5mm,bottom=3.5mm]
{\datafont\bfseries\fontsize{17pt}{17pt}\selectfont\textcolor{lm-accent}{#1}}\hspace{7pt}%
{\headingfont\bfseries\fontsize{11pt}{13pt}\selectfont\color{lm-black}#2}\par\vspace{3pt}
{\fontsize{9pt}{13pt}\selectfont\color{lm-secondary}#3}
\end{tcolorbox}}

\begin{document}
\makecover

\sectionlabel{Visión}
\section{La tesis}

\begin{callouttitled}{En una frase}
Makis transforma la URL de una empresa en un workspace de GTM investigado, documentado, visual y gobernable por una persona.
\end{callouttitled}

El problema no es que falten modelos que escriban. El problema es que el contexto del negocio está disperso: sitio web, competidores, estrategia, piezas y resultados viven en lugares distintos. El dueño termina copiando y pegando entre herramientas antes de poder decidir algo.

Makis resuelve tres cosas a la vez:
\begin{itemize}
  \item \textbf{Descubre:} usa la URL como la primera evidencia del negocio.
  \item \textbf{Estructura:} convierte evidencia en documentos, estrategia y piezas.
  \item \textbf{Gobierna:} deja al humano revisar, corregir y aprobar antes de ejecutar.
\end{itemize}

\begin{center}
\begin{tikzpicture}[node distance=5mm and 6mm, every node/.style={align=center}]
\node[rounded corners=3pt, fill=lm-black, text=white, minimum width=38mm, minimum height=17mm, font=\bfseries\small] (site) {Sitio web\\de la empresa};
\node[rounded corners=3pt, fill=lm-accent, text=white, minimum width=42mm, minimum height=17mm, font=\bfseries\small, right=of site] (makis) {Makis\\lee y organiza};
\node[rounded corners=3pt, fill=lm-accent-tint, draw=lm-accent, minimum width=42mm, minimum height=17mm, font=\bfseries\small, right=of makis] (workspace) {Workspace\\para operar};
\draw[-{Latex[length=2mm]}, thick, color=lm-accent] (site) -- (makis);
\draw[-{Latex[length=2mm]}, thick, color=lm-accent] (makis) -- (workspace);
\end{tikzpicture}
\end{center}

\section{El flujo de valor}

\begin{center}
\begin{tikzpicture}[node distance=2.5mm, every node/.style={align=center}]
\foreach \name/\label [count=\i] in {url/URL,investiga/Investiga,crea/Crea,revisa/Revisa,mide/Mide} {
  \node[rounded corners=2pt, fill=lm-surface, draw=lm-border, minimum width=26mm, minimum height=14mm, font=\bfseries\scriptsize] (\name) at ({(\i-1)*28mm},0) {\i. \\ \label};
}
\draw[-{Latex[length=1.8mm]}, thick, color=lm-accent] (url) -- (investiga);
\draw[-{Latex[length=1.8mm]}, thick, color=lm-accent] (investiga) -- (crea);
\draw[-{Latex[length=1.8mm]}, thick, color=lm-accent] (crea) -- (revisa);
\draw[-{Latex[length=1.8mm]}, thick, color=lm-accent] (revisa) -- (mide);
\draw[-{Latex[length=1.8mm]}, thick, color=lm-accent] (mide.south) to[out=-80,in=-100,looseness=.7] node[below,font=\scriptsize\color{lm-secondary}]{aprendizaje} (investiga.south);
\end{tikzpicture}
\end{center}

\sectionlabel{Sistema}
\section{Las cinco fases}

\phase{01}{Descubrimiento desde la URL}{Entrada: URL, objetivo y presupuesto. Makis extrae la oferta, categorías, tono y activos públicos. Salida: perfil del negocio y brief pendientes de confirmación.}
\phase{02}{Investigación y estrategia}{Exa y la lectura web encuentran competidores, tendencias y oportunidades. La estrategia fija audiencia, oferta, canales, mensajes y KPIs.}
\phase{03}{Plan y creatividades}{El creador produce prioridades, plan de contenido, copies, emails, anuncios, prompts, imágenes y una landing inicial.}
\phase{04}{Revisión y ejecución humana}{El operador edita en línea, regenera con feedback y usa controles explícitos: Aprobar, Editar o Rechazar. Resend puede ejecutar el correo; Ads se demuestran como mocks.}
\phase{05}{Resultados y aprendizaje}{Las métricas reales o simuladas producen leads, CPL, CAC y ROAS. El sistema propone el siguiente experimento y deja trazabilidad de la decisión.}

\section{El workspace es el producto}

\begin{center}
\begin{tikzpicture}[
item/.style={rounded corners=2pt, fill=lm-surface, draw=lm-border, minimum width=115mm, minimum height=8mm, align=left, font=\small}
]
\node[item] (a) {00-brief.md \hfill negocio, objetivo y supuestos confirmados};
\node[item, below=2mm of a] (b) {01-investigacion.md \hfill fuentes, competidores y hallazgos};
\node[item, below=2mm of b] (c) {02-estrategia.md \hfill audiencia, oferta, canales y KPIs};
\node[item, below=2mm of c] (d) {03-plan-de-contenido.md \hfill piezas, prioridades y calendario};
\node[item, below=2mm of d] (e) {04-creatividades.md \hfill copy, prompts y aprobación};
\node[item, below=2mm of e] (f) {05-resultados.md \hfill métricas, aprendizaje y próximo experimento};
\end{tikzpicture}
\end{center}

\sectionlabel{Experiencia}
\section{Una fuente de verdad, dos experiencias}

Markdown es durable, versionable, exportable y auditable. La UI no compite con él: lo convierte en la forma adecuada para decidir.

\begin{tabularx}{\textwidth}{>{\bfseries}p{25mm}p{40mm}X}
\toprule
& \textbf{Documento fuente} & \textbf{Vista en la interfaz} \\
\midrule
Contexto & Brief del negocio & Ficha para confirmar o corregir \\
Evidencia & Investigación & Tarjetas de competidores y fuentes \\
Dirección & Estrategia & Tablero de audiencia, oferta y KPIs \\
Producción & Creatividades & Studio para editar y aprobar \\
Aprendizaje & Resultados & Gráficos y próximo experimento \\
\bottomrule
\end{tabularx}

\vspace{12pt}
\begin{center}
\begin{tikzpicture}[node distance=12mm, every node/.style={align=center}]
\node[rounded corners=3pt, fill=lm-black, text=white, minimum width=44mm, minimum height=19mm, font=\bfseries\small] (md) {Markdown\\fuente de verdad};
\node[rounded corners=3pt, fill=lm-accent, text=white, minimum width=44mm, minimum height=19mm, font=\bfseries\small, right=of md] (ui) {UI renderizada\\para decidir};
\draw[-{Latex[length=2.5mm]}, thick, color=lm-accent] (md) -- node[above,font=\scriptsize\color{lm-secondary}]{renderiza} (ui);
\draw[-{Latex[length=2.5mm]}, thick, color=lm-accent] (ui.south) to[out=-100,in=-80] node[below,font=\scriptsize\color{lm-secondary}]{edita y versiona} (md.south);
\end{tikzpicture}
\end{center}

\section{Arquitectura y responsabilidades}

\begin{tabularx}{\textwidth}{>{\bfseries}p{46mm}X}
\toprule
\textbf{Capa} & \textbf{Responsabilidad} \\
\midrule
FastAPI + director & Orquestar onboarding y convertir URL en brief validable \\
Investigador + Exa & Extraer web, buscar mercado y guardar fuentes \\
Estratega + creador & Convertir hallazgos en dirección y activos de campaña \\
Workspace Markdown & Mantener documentos, versiones, decisiones y assets \\
SQLite & Indexar artefactos, aprobaciones, relaciones y eventos \\
UI + CopilotKit & Mostrar estado contextual y guiar decisiones \\
Analista + Chart.js & Medir, explicar y preparar la siguiente iteración \\
\bottomrule
\end{tabularx}

\vspace{12pt}
\begin{callouttitled}{Roles del equipo}
\textbf{Joel:} orquestación, SQLite, analítica y aprendizaje. \quad
\textbf{Diego:} URL, Exa e investigación. \quad
\textbf{Milu:} onboarding, cockpit, renderizado y revisión humana. \quad
\textbf{Freddy:} producto, demo y pitch.
\end{callouttitled}

\sectionlabel{Demo}
\section{Qué debe probar el hackathon}

\begin{callout}
\begin{itemize}
  \item Una URL real o preconfigurada crea un perfil y workspace en menos de un minuto.
  \item Al menos tres documentos Markdown aparecen como módulos de producto, no como texto plano.
  \item Una edición desde la interfaz actualiza el documento fuente y queda versionada.
  \item Una creatividad se aprueba y se ejecuta por Resend o mediante una simulación clara.
  \item El dashboard identifica un aprendizaje y lo conecta con el siguiente experimento.
\end{itemize}
\end{callout}

\begin{tcolorbox}[colback=lm-black,colframe=lm-black,arc=2mm,left=7mm,right=7mm,top=7mm,bottom=7mm]
{\headingfont\bfseries\fontsize{16pt}{22pt}\selectfont\color{white}
Makis no es una caja de chat. Entra por la URL de tu negocio, crea el sistema de trabajo y deja al humano gobernar cada decisión.}
\end{tcolorbox}

\end{document}
